\section{The Asymptotic Family}
\label{sec:scheme}

This section implements Section~\ref{sec:scheme} of the proof architecture.  There is one minor change in convention from this section onwards to keep in mind: any partition $F$ defined below corresponds to the partition $x\mapsto F(\sqrt2 x)$ in Definition~\ref{def:gaussian-correlation-function}.  This change of variables is harmless and keeps the Hermite covariance identity in its usual probabilist normalization.

\subsection{Hermite notation}

\begin{definition}[Orthonormal probabilist Hermite polynomials]\label{def:Hermite}
Let $(\He_n)_{n\ge0}$ denote the orthonormal probabilist Hermite polynomials.  Thus, for $Z\sim N(0,1)$,
\[
  \E[\He_m(Z)\He_n(Z)]=\delta_{mn}.
\]
In particular,
\[
  \He_0=1,
  \qquad \He_1(x)=x,
  \qquad \He_3(x)={x^3-3x\over\sqrt6},
  \qquad \He_5(x)={x^5-10x^3+15x\over\sqrt{120}}.
\]
If $(X,Y)$ are standard Gaussians with correlation $t$, then
\begin{equation}\label{eq:hermite-cov}
  \E[\He_m(X)\He_n(Y)]=\delta_{mn}t^n.
\end{equation}
We use this as our normalization throughout the coefficient computation.  Standard references for these identities include~\cite[Sec.~6.1]{AAR99}.
\end{definition}

\begin{definition}[Gaussianized Hermite-chaos average]\label{def:gaussianized-chaos}
Fix an odd integer $d\ge3$.  A length-$N$ \emph{Gaussianized degree-$d$ Hermite-chaos average} is the random variable
\[
  P_{d,N}=N^{-1/2}\sum_{j=1}^N \He_d(U_{d,j}),
\]
where the $U_{d,j}$ are independent standard Gaussians.  The word ``Gaussianized'' refers to the fact that, for fixed $d$, the variable $P_{d,N}$ converges in distribution to a Gaussian as $N\to\infty$, while paired copies retain the covariance power of the degree-$d$ Wiener chaos.  Specifically, if $(U_{d,j},U_{d,j}')$ are paired with correlation $t$, then \eqref{eq:hermite-cov} gives
\[
  \E[P_{d,N}P_{d,N}']=t^d.
\]
For background on Hermite polynomials and Wiener chaos, see Janson's book on Gaussian Hilbert spaces~\cite{Janson1997GaussianHilbertSpaces}.
\end{definition}

\subsection{Finite-dimensional schemes}\label{subsec:finite-schemes}

\begin{definition}[Finite Hermite-average family]\label{def:finite-family}
Let $D\subseteq\{3,5,7,9,\ldots\}$ be a finite set of odd degrees.  For each $d\in D$, fix a real weight $s_d$, and define the relative sign
\[
  \sigma_d=(-1)^{(d-1)/2}.
\]
% Thus $\sigma_3=-1$, $\sigma_5=+1$, $\sigma_7=-1$, and so on.

For length $N$, let
\[
  P_{d,N}=N^{-1/2}\sum_{j=1}^N \He_d(U_{d,j}),
\]
where all variables $Z,X,U_{d,j}$ are independent standard Gaussians.  The finite rounding functions are
\begin{align}
  f_N&=\sgn\biggl(Z+\eta \He_3(X)+\sum_{d\in D}s_dP_{d,N}\biggr),\label{eq:finite-f-general}\\
  g_N&=\sgn\biggl(Z-\eta \He_3(X)+\sum_{d\in D}\sigma_d s_dP_{d,N}\biggr).\label{eq:finite-g-general}
\end{align}
They are odd measurable functions on $\R^{2+|D|N}$.  
% The value of $\sgn(0)$ is immaterial, since every threshold event has probability zero.
\end{definition}

\begin{definition}[The third/fifth finite scheme]\label{def:finite-35}
The scheme used for the numerical certificate is the specialization of Definition~\ref{def:finite-family} with $D=\{3,5\}$ and
\begin{equation}\label{eq:params}
  \eta=0.136419125,
  \qquad
  s_3=0.34101124,
  \qquad
  s_5=0.05276111.
\end{equation}
Equivalently,
\begin{align}
  f_N&=\sgn\bigl(Z+\eta \He_3(X)+s_3P_{3,N}+s_5P_{5,N}\bigr),\label{eq:finite-f-35}\\
  g_N&=\sgn\bigl(Z-\eta \He_3(X)-s_3P_{3,N}+s_5P_{5,N}\bigr).\label{eq:finite-g-35}
\end{align}
Thus degree $3$ is anti-aligned between the two sides, while degree $5$ is aligned.
\end{definition}

\subsection{The limiting model}\label{subsec:limiting-model}

\begin{definition}[Limiting model]\label{def:limiting-model}
Let $D,\eta,(s_d)_{d\in D}$ be as in Definition~\ref{def:finite-family}.  Let $R_d$, $d\in D$, be independent standard Gaussians.  The limiting threshold pair is
\begin{align}
  f(z,x,(r_d))&=\sgn\biggl(z+\eta \He_3(x)+\sum_{d\in D}s_dr_d\biggr),\label{eq:lim-f-general}\\
  g(z,x,(r_d))&=\sgn\biggl(z-\eta \He_3(x)+\sum_{d\in D}\sigma_d s_dr_d\biggr).\label{eq:lim-g-general}
\end{align}
For the limiting correlation at parameter $t$, the paired coordinates satisfy
\[
  \rho_Z=t,
  \qquad \rho_X=t,
  \qquad \rho_d=t^d\quad(d\in D),
\]
and different coordinate pairs are independent.  The arcsine-normalized limiting correlation is
\begin{equation}\label{eq:lim-H}
  H(t)={\pi\over2}\E\bigl[f(Z,X,(R_d))g(Z',X',(R_d'))\bigr].
\end{equation}
\end{definition}

Because all degrees in the model are odd, $H$ is an odd function.  For $D=\{3,5\}$, Definition~\ref{def:limiting-model} implies that the extra coordinates have correlations $t^3$ and $t^5$, which come from the finite Hermite averages through \eqref{eq:hermite-cov}.
